% Body of the theoretical appendix; \input by main.tex.
\noindent
This appendix states the model in the form used by the algorithm, proves that the
algorithm returns the exact Pareto front, and derives its complexity. Every result
is linked to the experiment that checks it numerically (\cref{app:verification}).
Throughout, \emph{dominance} means Pareto dominance: an outcome is weakly better in
every objective and strictly better in at least one. Two plans with identical
objective vectors are treated as one point of the front.

%=====================================================================================
\subsection{Model and notation}\label{app:model}

The planning horizon has periods $t=1,\dots,T$ and a deterministic lead time $L$, so
orders can be placed in the decision periods $\tau=1,\dots,\Tp$ with $\Tp=T-L$; an order
placed in period $\tau$ is received in period $\tau+L$. Demand $D(t)\ge 0$ is given
(one sample path of the stochastic demand), $\cumD(t)=\sum_{s\le t}D(s)$ with
$\cumD(0)=0$, and $I_0$ is the initial stock. Order quantities lie in
\begin{equation}
  \cQ=\{0\}\cup\{\moq+a\,r : a\in\mathbb{N}_0\},
\end{equation}
where $\moq$ is the minimum order quantity and $r$ the rounding value. The outer
(policy) variables are the safety stock $\Ss$ and the safety order time $\SOT$, taken
from the domain
\begin{equation}
  \Omega=\{(\Ss,\SOT)\in\mathbb{R}^2_{\ge 0} : \SOT+\Ss/\mu\le\MS\},
\end{equation}
or from a finite grid in $\Omega$ (\emph{grid mode}).

\paragraph{Cumulative reformulation.}
Let $X(\tau)=\sum_{s\le\tau}Q(s)$, $X(0)=0$. Then the stock at the end of period
$\tau+L$ is $I(\tau+L)=I_0+X(\tau)-\cumD(\tau+L)$, while $I(t)=I_0-\cumD(t)$ for
$t\le L$ does not depend on the decisions. The service requirement
$I(t)\ge\Ss+\SOT\cdot D(t)$ for $t=L+1,\dots,T$ becomes
\begin{equation}\label{eq:req}
  X(\tau)\ge\beta_\tau(\Ss,\SOT)
  :=\Ss+\SOT\cdot D(\tau+L)+\cumD(\tau+L)-I_0,\qquad \tau=1,\dots,\Tp,
\end{equation}
and the average stock $\Ibar(t)=\tfrac12\bigl(I(t-1)+I(t)\bigr)$ in period $t=\tau+L$ is
\begin{equation}\label{eq:ibar}
  \Ibar(\tau+L)=b_\tau+\tfrac12\bigl(X(\tau-1)+X(\tau)\bigr),\qquad
  b_\tau:=I_0-\tfrac12\bigl(\cumD(\tau+L-1)+\cumD(\tau+L)\bigr).
\end{equation}
Inventory turnover is defined only for positive average stock, so we impose
$\Ibar(t)\ge\eps$ for all $t$, with a fixed threshold $\eps>0$ (the implementation
uses $\eps=1$ unit).

\paragraph{Objectives.}
With $k=\#\{\tau : Q(\tau)>0\}$ the number of orders, the four objectives of
Kania et al.\ read
\begin{align}
  \POC &= p\,X(\Tp)+c\,k &&\text{(min)},\label{eq:poc}\\
  \HC  &= h\sum_{t=1}^{T}\Ibar(t)
        = h\Bigl[C_0+\sum_{\tau=1}^{\Tp}X(\tau)-\tfrac12X(\Tp)\Bigr]
        &&\text{(min)},\label{eq:hc}\\
  \CSL &= \Phi\bigl(B/\sigma\bigr),\quad B:=\Ss+\mu\,\SOT &&\text{(max)},\label{eq:csl}\\
  \ITO &= A_0+\sum_{\tau=1}^{\Tp}\frac{D(\tau+L)+\sigma}{\Ibar(\tau+L)} &&\text{(max)},\label{eq:ito}
\end{align}
where $C_0=\sum_{t\le L}\Ibar(t)+\sum_{\tau}b_\tau$ and
$A_0=\sum_{t\le L}(D(t)+\sigma)/\Ibar(t)$ are constants. The second form of
\eqref{eq:hc} follows from \eqref{eq:ibar} and $X(0)=0$. We write
$F=(\POC,\HC,-\CSL,-\ITO)$, to be minimised componentwise.

\paragraph{Assumptions.} Both instances studied satisfy the following.
\begin{assumption}\label{as:initial}
  $I(L)=I_0-\cumD(L)\ge 0$ and $\Ibar(t)\ge\eps$ for $t\le L$.
\end{assumption}
\begin{assumption}\label{as:lattice}
  $\moq$ is a positive multiple of $r$, i.e.\ $\moq=m r$ with $m\in\mathbb{N}$.
\end{assumption}
Under \cref{as:lattice}, $\cQ\setminus\{0\}=\{r j : j\ge m\}$ is closed under addition,
and the set of reachable cumulative quantities is
$\cL=\{0\}\cup\{r j : j\ge m\}\subset r\mathbb{N}_0$. For $y\in\mathbb{R}$ let
$\lceilL{y}=\min\{\ell\in\cL : \ell\ge y\}$. Because $X(\tau)\in\cL$,
\begin{equation}\label{eq:kappa}
  X(\tau)\ge\beta_\tau \iff X(\tau)\ge\kappa_\tau,\qquad
  \kappa_\tau:=\lceilL{\beta_\tau}.
\end{equation}
We call $\kappa=(\kappa_1,\dots,\kappa_{\Tp})$ the \emph{rounded requirement vector}.

%=====================================================================================
\subsection{The inner problem as a multi-criteria shortest path}\label{app:inner}

Fix an outer point $z=(\Ss,\SOT)$. Then $\CSL$ is constant, and a plan is a
non-decreasing path $X(0)=0,X(1),\dots,X(\Tp)$ in $\cL$ whose increments lie in $\cQ$.
Consider the layered directed acyclic graph whose nodes are pairs $(\tau,x)$ with
$x\in\cL$, and whose arcs $(\tau-1,x')\to(\tau,x)$ exist when
\begin{equation}\label{eq:arc}
  x-x'\in\cQ,\qquad x\ge\kappa_\tau,\qquad b_\tau+\tfrac12(x'+x)\ge\eps .
\end{equation}
Each arc carries the additive cost vector
\begin{equation}\label{eq:arccost}
  \bigl(\,\mathbb{1}[x>x'],\; x,\; \tfrac{D(\tau+L)+\sigma}{b_\tau+\frac12(x'+x)}\,\bigr),
\end{equation}
and a path to $(\tau,x)$ accumulates the label $\ell=(k,W,R)$: the number of orders, the
partial sum $W=\sum_{s\le\tau}X(s)$, and the partial turnover $R$. At a terminal node
$(\Tp,x)$, \eqref{eq:poc}--\eqref{eq:ito} give
\begin{equation}\label{eq:terminal}
  \POC=p x+c k,\qquad \HC=h\bigl(C_0+W-\tfrac12x\bigr),\qquad \ITO=A_0+R .
\end{equation}
For a fixed terminal node these are strictly monotone in the label (increasing in $k$
and $W$, increasing in $R$ for the maximised $\ITO$). Hence two paths ending at the same
node compare in $F$ exactly as their labels compare under the order
\begin{equation}
  (k,W,R)\preceq(k',W',R')\iff k\le k',\ W\le W',\ R\ge R'.
\end{equation}

The label-setting algorithm processes the layers $\tau=1,\dots,\Tp$ in order, extends
every stored label along every outgoing arc, and keeps at each node only the labels that
are non-dominated under $\preceq$. The implementation stores them as one
two-dimensional staircase in $(W,R)$ per value of $k$; this is a data structure for the
same dominance test and does not affect correctness.

%=====================================================================================
\subsection{Reductions of the search space}\label{app:reductions}

The algorithm searches a restricted graph: states are bounded above (\cref{lem:ub}), and
ordering arcs that can never be Pareto-relevant are removed (\cref{lem:D}). Both lemmas
are proved by exhibiting, for every plan that the restriction excludes, a feasible plan
with the same outer point that dominates it. So no Pareto-optimal plan is excluded.

\begin{lemma}[Upper bound on the cumulative quantity]\label{lem:ub}
  Let $\hat\beta=\max\{0,\max_\tau\beta_\tau\}$ and let \cref{as:initial} hold. Every
  feasible plan with $X(\Tp)\ge\hat\beta+\max\{\moq,r\}+2\eps$ is dominated by a feasible
  plan with the same outer point. Consequently every Pareto-optimal plan satisfies
  $X(\Tp)<\hat\beta+\max\{\moq,r\}+2\eps$.
\end{lemma}

\begin{proof}
  Let $\tau^\ast$ be the last period with $Q(\tau^\ast)>0$ (it exists because
  $X(\Tp)>0$), so $X(\tau)=X(\Tp)$ for $\tau\ge\tau^\ast$. Reduce the last order by
  $\delta=r$ if $Q(\tau^\ast)\ge\moq+r$, and by $\delta=\moq$ (that is, remove it) if
  $Q(\tau^\ast)=\moq$. In both cases $Q(\tau^\ast)-\delta\in\cQ$, $\delta\le\max\{\moq,r\}$,
  and the new plan $X'$ equals $X$ before $\tau^\ast$ and $X-\delta$ from $\tau^\ast$ on;
  it remains non-decreasing.

  \emph{Feasibility.} For $\tau\ge\tau^\ast$,
  $X'(\tau)=X(\Tp)-\delta\ge\hat\beta+2\eps\ge\beta_\tau+2\eps$, so \eqref{eq:req} holds and
  $I'(\tau+L)=X'(\tau)-\beta_\tau+\Ss+\SOT\,D(\tau+L)\ge 2\eps$. Hence
  $\Ibar'(t)\ge2\eps$ for $t>\tau^\ast+L$. In period $t=\tau^\ast+L$,
  $\Ibar'(t)=\tfrac12\bigl(I(t-1)+I'(t)\bigr)\ge\tfrac12(0+2\eps)=\eps$, because $I(t-1)$
  is unchanged and non-negative: by \cref{as:initial} if $\tau^\ast=1$, and by
  \eqref{eq:req} in period $\tau^\ast-1$ otherwise. Earlier periods are unchanged.

  \emph{Dominance.} $\POC$ decreases by $p\delta>0$ (and by a further $c$ if the order is
  removed). By \eqref{eq:hc}, $\HC$ changes by
  $h\bigl[-\delta(\Tp-\tau^\ast+1)+\delta/2\bigr]<0$. $\CSL$ is unchanged. Every
  $\Ibar(t)$ either stays the same or decreases, so $\ITO$ does not decrease.
\end{proof}

The implementation therefore uses the state set
$\{x\in\cL : x\le\lceil\hat\beta\rceil+\max\{\moq,r\}+\lceil2\eps\rceil\}$, where
$\hat\beta$ is the maximum of $\beta_\tau$ over the whole outer domain. This is one
superset that serves every outer point.

\begin{lemma}[Order only when needed]\label{lem:D}
  Let \cref{as:lattice} hold, and consider a feasible plan with an order $q=Q(\tau)>0$ and
  $x'=X(\tau-1)$. Suppose that
  \begin{enumerate}[label=(\alph*),nosep]
    \item $x'\ge\beta_\tau$ (equivalently $x'\ge\kappa_\tau$),
    \item $b_\tau+x'\ge\eps$, and
    \item $\tau=\Tp$, or $b_{\tau+1}+x'+q/2\ge\eps$.
  \end{enumerate}
  Then the plan is dominated by a feasible plan with the same outer point.
\end{lemma}

\begin{proof}
  If $\tau<\Tp$, move the order to the next period:
  $Q'(\tau)=0$ and $Q'(\tau+1)=Q(\tau+1)+q$. This quantity lies in $\cQ$, either because
  $Q(\tau+1)=0$ or because $\cQ\setminus\{0\}$ is closed under addition
  (\cref{as:lattice}). If $\tau=\Tp$, drop the order. The new plan differs from the old
  one only in $X'(\tau)=x'$.

  \emph{Feasibility.} The requirement in period $\tau$ holds by (a), and all other
  requirements are unchanged. By \eqref{eq:ibar} and (b),
  $\Ibar'(\tau+L)=b_\tau+x'\ge\eps$. If $\tau<\Tp$, then $X(\tau+1)\ge X(\tau)=x'+q$, and (c)
  gives $\Ibar'(\tau+L+1)=b_{\tau+1}+\tfrac12\bigl(x'+X(\tau+1)\bigr)\ge b_{\tau+1}+x'+q/2\ge\eps$.
  No other period changes.

  \emph{Dominance.} If $\tau<\Tp$, then $X(\Tp)$ is unchanged and $k$ stays the same or
  decreases by one, so $\POC$ does not increase. $\sum_\tau X(\tau)$ decreases by $q$, so
  $\HC$ decreases strictly. If $\tau=\Tp$, then $\POC$ decreases and $\HC$ changes by
  $h(-q+q/2)<0$. $\CSL$ is unchanged, and $\ITO$ does not decrease because the average
  stock in periods $\tau+L$ and $\tau+L+1$ does not increase.
\end{proof}

\begin{remark}[Why condition (c) is needed]\label{rem:c}
  Postponing an order lowers the average stock in \emph{two} periods, $\tau+L$ and
  $\tau+L+1$. A version of \cref{lem:D} that checks only (a) and (b) is not valid: the
  postponed plan may violate $\Ibar\ge\eps$ in period $\tau+L+1$, so the claimed dominating
  plan need not exist. In our experiments the earlier implementation, which omitted (c),
  missed 2 of 19\,372 Pareto points in one instance (SKU~2, seed~42, grid mode; both
  points have $\Ss=\SOT=0$, where the stock is lowest). The corrected rule reproduces the
  front of the DP without any pruning. All 20 instances of experiment E2 (seeds 1--10) are
  unaffected.
\end{remark}

\begin{remark}[Implementation]
  For fixed $(\tau,x')$, condition (c) is monotone in $q$, and (a) and (b) do not depend
  on $q$. Once an order size is prunable, every larger one is too. The implementation
  enumerates target states in increasing order and stops at the first prunable order.
\end{remark}

%=====================================================================================
\subsection{Exactness of the algorithm}\label{app:exact}

\begin{theorem}[Exactness for a fixed outer point]\label{thm:exact}
  Let \cref{as:initial,as:lattice} hold and fix $z\in\Omega$. Let $G_z$ be the graph of
  \cref{app:inner} restricted to the state set of \cref{lem:ub}, with the arcs excluded by
  \cref{lem:D} removed. The label-setting algorithm on $G_z$, followed by evaluating
  \eqref{eq:terminal} and filtering non-dominated outcomes over all terminal nodes,
  returns exactly the set of non-dominated outcome vectors of the inner problem at $z$.
\end{theorem}

\begin{proof}
  \emph{(i) The restriction loses no Pareto-optimal plan.} A Pareto-optimal plan is not
  dominated by any feasible plan. By \cref{lem:ub} it lies within the state bound, and by
  \cref{lem:D} none of its arcs is pruned. So it is a path of $G_z$. Conversely, every
  path of $G_z$ satisfies \eqref{eq:arc} and is therefore a feasible plan.

  \emph{(ii) Label dominance is preserved by extension.} Arc existence and arc cost
  depend only on the arc $(\tau-1,x')\to(\tau,x)$, not on the path that reached
  $(\tau-1,x')$. This holds for the pruning rule as well, which depends on
  $(\tau,x',x-x')$ only. If labels $\ell\preceq\ell'$ sit at the same node, extending both
  along the same arc preserves $\preceq$, because the costs \eqref{eq:arccost} add
  componentwise.

  \emph{(iii) Induction over layers.} Suppose that at layer $\tau-1$ every path of $G_z$
  to a node has a stored label at that node which weakly dominates it under $\preceq$.
  Extending that stored label along the path's next arc weakly dominates the extended
  path, by (ii). Either it is kept, or it is discarded because another kept label
  dominates it; by transitivity, some kept label dominates the extended path. So the claim
  holds at layer $\tau$, and it holds trivially at layer $0$.

  \emph{(iv) Terminal evaluation.} By (iii) and the monotonicity of \eqref{eq:terminal},
  every plan in $G_z$ ending at $(\Tp,x)$ is weakly dominated in $F$ by a returned outcome
  at the same node. Combining this with (i), the returned set contains every
  non-dominated outcome. It contains only feasible outcomes, so after the final filter it
  equals the set of non-dominated outcome vectors.
\end{proof}

\begin{lemma}[The inner problem depends on $z$ only through $\kappa$]\label{lem:memo}
  For $z,z'\in\Omega$ with $\kappa(z)=\kappa(z')$, the graphs $G_z$ and $G_{z'}$ coincide,
  so the inner problems have the same outcomes in $(\POC,\HC,\ITO)$. If
  $\kappa(z)\le\kappa(z')$ componentwise, every plan feasible at $z'$ is feasible at $z$.
\end{lemma}

\begin{proof}
  By \eqref{eq:kappa}, $z$ enters \eqref{eq:arc} only through $\kappa_\tau$. Condition
  (a) of \cref{lem:D} is equivalent to $x'\ge\kappa_\tau$, and conditions (b) and (c) do
  not involve $z$. The objectives \eqref{eq:terminal} do not involve $z$, and the state
  bound is common to all outer points. The second statement follows because
  $x\ge\kappa_\tau(z')\ge\kappa_\tau(z)$.
\end{proof}

\begin{corollary}[Exact memoization]\label{cor:memo}
  Among outer points with the same rounded requirement vector, only one with the largest
  $B=\Ss+\mu\SOT$ can contribute to the Pareto front: the others have identical
  $(\POC,\HC,\ITO)$ outcomes and a strictly smaller $\CSL$, because $\Phi$ is strictly
  increasing. The algorithm therefore solves one inner problem per distinct vector
  $\kappa$, attaches the largest $B$ of its class, and applies a final four-objective
  filter over the classes. By \cref{thm:exact} the result is the exact front over the
  outer candidate set.
\end{corollary}

%=====================================================================================
\subsection{The continuous outer domain}\label{app:continuous}

Each $\beta_\tau$ is affine in $z$: $\beta_\tau(z)=\Ss+D_\tau\SOT+c_\tau$ with
$D_\tau=D(\tau+L)$ and $c_\tau=\cumD(\tau+L)-I_0$. Hence $\kappa_\tau(z)$ is piecewise
constant. It can change only on the lines
\begin{equation}
  H_{\tau,\lambda}=\{z : \Ss+D_\tau\SOT=\lambda-c_\tau\},\qquad
  \lambda\in\Lambda:=\cL\cap\bigl[0,\lceilL{\hat\beta}\bigr],
\end{equation}
where $\hat\beta$ is the maximum of $\beta_\tau$ over $\Omega$. Let $\mathcal{A}$ be the
arrangement formed by these lines together with the three lines that bound $\Omega$
($\Ss=0$, $\SOT=0$, $\SOT+\Ss/\mu=\MS$). Let $V$ be the finite set of its vertices,
that is, the intersection points of two non-parallel lines of $\mathcal{A}$ that lie in
$\Omega$.

\begin{theorem}[Vertex sufficiency]\label{thm:vertex}
  For every $z\in\Omega$ there is a vertex $v\in V$ with $\kappa(v)\le\kappa(z)$
  componentwise and $B(v)\ge B(z)$. Consequently the set of non-dominated outcome vectors
  over all of $\Omega$ equals the set of non-dominated outcome vectors over $V$.
\end{theorem}

\begin{proof}
  Let $P(z)=\{w\in\Omega : \beta_\tau(w)\le\kappa_\tau(z)\ \forall\tau\}$. It is the
  intersection of the compact polygon $\Omega$ with finitely many closed half-planes. It
  is therefore a compact convex polygon, and it contains $z$ because
  $\beta_\tau(z)\le\kappa_\tau(z)$. The linear function $B$ attains its maximum over
  $P(z)$ at a vertex $v$ of $P(z)$, so $B(v)\ge B(z)$. Every $w\in P(z)$ satisfies
  $\kappa_\tau(w)\le\kappa_\tau(z)$, because $\kappa_\tau(z)\in\cL$ is an upper bound of
  $\beta_\tau(w)$ and $\kappa_\tau(w)$ is the smallest such element of $\cL$. A vertex of
  $P(z)$ is the intersection of two non-parallel boundary lines of $P(z)$. Each such line
  is either a boundary line of $\Omega$ or a line $\beta_\tau=\kappa_\tau(z)$, that is,
  $H_{\tau,\lambda}$ with $\lambda=\kappa_\tau(z)\in\Lambda$. Hence $v\in V$.

  For the consequence, let $y=F(z,X)$ be non-dominated over $\Omega$. By
  \cref{lem:memo}, the plan $X$ is feasible at $v$, and $F(v,X)$ has the same
  $(\POC,\HC,\ITO)$ and $\CSL(v)\ge\CSL(z)$. So $F(v,X)$ weakly dominates $y$. Since $y$
  is non-dominated, $F(v,X)=y$; thus $y$ is attained at a vertex. The reverse inclusion
  holds because $V\subseteq\Omega$.
\end{proof}

\begin{remark}[Correction of an earlier claim]
  An earlier version of this work asserted that the continuous front need not be closed,
  so that some cells of the arrangement contributed points only as a supremum of
  $\CSL$. To catch these, it added one interior probe per angular sector around each
  vertex. \Cref{thm:vertex} shows the probes are unnecessary. The vertex of a cell where
  the supremum is attained belongs to a neighbouring cell with a componentwise smaller
  requirement vector. Its outcome reaches the supremum and weakly dominates every point of
  the open cell. Numerically, in three instances (SKU~1 seed~1; SKU~2 seeds~3 and~6),
  vertices alone and vertices plus probes give identical objective sets (52\,715,
  46\,875 and 67\,989 points). All points that the probe version had marked as
  ``not attained'' are duplicates of attained points.
\end{remark}

By \cref{cor:memo,thm:vertex}, the continuous front is obtained by evaluating $\kappa$ at
the vertices, keeping the largest $B$ per distinct vector, solving one inner problem per
vector, and filtering. In our instances, $\lvert V\rvert$ is several hundred to a few
thousand at $T=48$, and the number of distinct vectors is 7--26\% smaller.

%=====================================================================================
\subsection{Invariance with respect to the ordering cost}\label{app:invariance}

\begin{proposition}[Ordering-cost invariance]\label{prop:c}
  Let \cref{as:lattice} hold and $0<c<p\,r/\Tp$. Then the dominance relation between any
  two candidate solutions (outer point and plan) does not depend on $c$. In particular,
  the Pareto-optimal set of decisions is the same for every $c$ in this interval.
\end{proposition}

\begin{proof}
  $c$ enters only through $\POC=pX(\Tp)+ck$, so it suffices to show that the sign of
  $\Delta\POC=p\,\Delta X+c\,\Delta k$ between two solutions does not depend on $c$. By
  \cref{as:lattice}, $X(\Tp)\in r\mathbb{N}_0$. If $\Delta X\ne0$, then
  $\lvert\Delta X\rvert\ge r$ and $\lvert\Delta k\rvert\le\Tp$, so
  $p\lvert\Delta X\rvert\ge p r>c\Tp\ge c\lvert\Delta k\rvert$, and $\Delta\POC$ has the
  sign of $\Delta X$. If $\Delta X=0$, then $\Delta\POC=c\,\Delta k$ has the sign of
  $\Delta k$ because $c>0$. In both cases the sign, including the case
  $\Delta\POC=0$, is determined independently of $c$.
\end{proof}

\begin{remark}
  For SKU~1 ($p=134$, $r=70$) at $T=48$ ($\Tp=43$) the interval is $0<c<218.1$, which
  contains the values $c=50$ and $c=200$ used in experiment E4. For $c=800$ we also
  observed identical decision sets, which shows that the condition is sufficient but not
  necessary. For SKU~2 ($r=3$) the interval is narrow ($c<33$), and the proposition makes
  no claim at the nominal $c=200$.
\end{remark}

%=====================================================================================
\subsection{Complexity}\label{app:complexity}

Let $S$ be the number of states per layer. With $\hat\beta\approx\mu T$ and spacing $r$
in $\cL$, $S=\Theta(\mu T/r)$. Under \cref{as:lattice}, the partial sum $W$ is a multiple
of $r$ not exceeding $\Tp x_{\max}$, where $x_{\max}$ is the state bound of \cref{lem:ub};
it therefore takes at most $\Tp S$ values, and a node holds labels for at most
$(\Tp+1)\,\Tp S$ distinct pairs $(k,W)$. For each pair only the largest $R$ survives, so
the number $N$ of labels per node satisfies
\begin{equation}
  N=O(\Tp^{2}S).
\end{equation}
A layer has $O(S^2)$ arcs, and each arc extends at most $N$ labels. One inner
problem therefore costs
\begin{equation}\label{eq:complexity}
  O\bigl(\Tp\,S^{2}\,N\log N\bigr)
  =O\bigl(T^{3}S^{3}\log(TS)\bigr)
  =O\bigl(T^{6}(\mu/r)^{3}\log(T\mu/r)\bigr),
\end{equation}
where the logarithm comes from staircase insertion. This is a pseudo-polynomial bound. It
is polynomial in $T$ for a fixed ratio $\mu/r$. The empirical exponents of experiment E3
agree with it: time grows as $T^{6.0\text{--}6.4}$ and as $(\mu/r)^{3.5}$, the stored
labels grow as $T^{4.9\text{--}5.0}$, and time grows as $(\text{labels})^{1.25}$. The
bound is not tight in practice, since \cref{lem:D} removes most ordering arcs and the
staircases remain far below $N$. The number of inner problems is the number of
distinct vectors $\kappa$, which is at most $\lvert V\rvert=O\bigl((\Tp\lvert\Lambda\rvert)^2\bigr)$
for the continuous domain (here $\lvert\Lambda\rvert$ is the number of relevant lattice
values) and at most the grid size in grid mode.

%=====================================================================================
\subsection{Numerical verification of the theoretical results}\label{app:verification}

\Cref{tab:verification} lists, for each result, the experiment that checks it with the
corrected code. Experiment labels refer to \cref{sec:design}.

\begin{table}[h]
  \centering
  \caption{Results of this appendix and their numerical verification.}
  \label{tab:verification}
  \small
  \begin{tabular}{@{}p{3.6cm}p{11.2cm}@{}}
    \toprule
    Result & Verification \\
    \midrule
    \Cref{thm:exact} with \cref{lem:ub,lem:D}
      & E1: DP equals exhaustive enumeration (state bound three times larger, no pruning)
        on 100/100 random instances; at $T=48$ (SKU~2, seed~42) the DP with pruning equals
        the DP without pruning (\cref{rem:c}); all 40 E2 fronts (2 SKUs $\times$ 10 seeds
        $\times$ grid/continuous) are identical before and after the correction. E5a: on 10 + 16 outer points at $T=48$, the
        $\varepsilon$-constraint MILP solved with HiGHS gives identical fronts. \\
    \Cref{thm:vertex}
      & E1: for 20 instances (2 SKUs $\times$ 10 seeds) $\times$ $10^6$ random points of
        $\Omega$ (one quarter on the upper boundary), every point has a vertex with
        $\kappa(v)\le\kappa(z)$ and $B(v)\ge B(z)$; front with and without sector probes
        identical (3 instances); all 20 continuous E2 fronts unchanged after removing the
        probes. \\
    \Cref{prop:c}
      & E4: identical decision sets for $c\in\{50,200\}$ (covered) and $c=800$ (beyond
        the sufficient condition). \\
    Exactness overall
      & E5b: none of the exact points is dominated by any of about $1.2\times10^8$
        NSGA-III evaluations. \\
    Complexity \eqref{eq:complexity}
      & E3: fitted exponents reported above. \\
    \bottomrule
  \end{tabular}
\end{table}
